\chapter*{Preface}

\vspace{\baselineskip}
\lettrine{W}{hile} the main characters in this book are entirely fictitious, their idle musings, spirited debates, and earnest explorations into the numbered world are inspired by real mathematical and physical concepts. These characters draw upon the rich tapestry of past discoveries, referencing the giants upon whose shoulders modern mathematics and physics stand. Figures like Nikolai Lobachevsky, Emmy Noether, Srinivasa Ramanujan, and Claude Shannon make their appearances not in person, but woven into their dialogues and contemplations.

This is a journey into the interconnected realms of number theory, geometry, and physics. The characters' dialogue often begin on firm, well-trodden turf, grounded in established theories. Yet, they are encouraged to wander in speculating, hypothesizing, and leaping beyond established rigor and dogma. Their flights of fancy have a more computational tone, hinting to the younger mathematician that there are as many \textit{mule} as their are\textit{scramble} paths to the summit of the forbidding mathematical edifice. Their musings lay no claim to the authority or rigour of the prodigious formal workings of a professional mathematician. They merely serve as an invitation to the reader to embrace that playful curiosity that accompanies genuine discovery.

% \subsection*{Fiction as a Gateway to Reality}
The interactions and musings of these fictional characters are designed to bring otherwise abstract ideas somewhat to life. References to Lobachevsky’s pioneering work on hyperbolic geometry and Shannon’s insights into information theory serve as genuine contributions, with the characters’ debates and pontifications designed to make such complex ideas relatable and prematurely accessible. In truth, they may at best serve as a primer: as a means of deceiving the mind into a false sense of familiarity, to be exposed in earnest later in one’s studies.\footnote{As John von Neumann is reputed to have said: “In mathematics, you don’t understand things. You just get used to them.”} The hope is that their idlings may occasionally echo your own, and that their speculative detours might spark connections you hadn’t previously considered. This is barely a novel, and no more a text book offering little in the way of rigor for any claims. Instead, it offers an insight into computational and differential geometric number theory alluding to some of their cryptographic uses within the financial sector. You the youthful reader, are invited to disentangle the speculative flights of fancy from the well grounded, useful insights provided. Somewhat mirroring the essence of mathematical inquiry the reader might want to run with a couple of ideas building on the referenced python code, or indeed more through formal mathematical lines of attack.

Mathematics, at its core, can be both a deadly serious discipline and a playful art. Here we celebrate this duality demonstrating some of the joy of unpicking the chiffon fabric that is this numbered universe of ours.

\subsection*{Acknowledgments and Disclaimer}
While care has been taken to ensure allusions to past mathematicians and their contributions are apt, the speculative extensions and interpretations of their work are purely fictional. Any resemblance of the fictional characters to real persons, living or deceased, is entirely coincidental.

This book is dedicated to the spirit of the youth. To those among them who, despite wading through the relics of the ancients and the laboured algebra of Enlightenment minds, are revealed as the freshest chorizos from the great sausage factory of our schooling, still delighting in the numbers that thread the universe together.
\begin{flushright}
\textit{— L McCulloch-James}
\end{flushright}

%%%%%%%%%%%%%%%%%%%%%%

\section*{Plot Overview for those losing the thread}

Eleanor Price, the intellectual matriarch of the Merriweather Circle, has long envisioned an independent mathematical institute housed within the Merriweather Library. However, financial constraints threaten this vision.

Demi-Moore, a brilliant mathematician with ties to biomedical sciences, sees the potential for interdisciplinary research and proposes an investment opportunity—one that would provide the much-needed financial backing for the institute. However, this offer comes with an implicit compromise.

Eleanor Price, Merriweather Circle guiding intellectual force, has long resisted financial entanglements, believing in the sanctity of pure mathematics. But reality looms: without funding, the Circle may cease to exist.

\subsection*{1. Demi-Moore’s Industry Play}

Demi-Moore, once deeply embedded in the world of \textbf{big pharma}, retains connections that few within the Circle would dare engage with. Her former partner, \textbf{Joe ("GI")}, leads a research-driven investment firm specializing in the \textit{mathematics of biological systems}. Celiac Wu, now debt-laden drawn by the possibility of securing research space, approaches Eleanor with a proposition.

\smallskip

\textbf{The pitch:} A partnership between the Merriweather Circle and biomedical investors, using mathematical modeling to explore \textit{metabolic scaling laws in longevity research}.  

\smallskip

But the deal comes with consequences. <<Would the Circle become beholden to corporate interests? Would the research be steered toward commercial applications rather than fundamental understanding?>>

\subsection*{2. Eleanor's Moral Dilemma}

Eleanor hesitates. She has always believed that the pursuit of truth should not be dictated by profit motives. But without intervention, the library—and the Circle—could be lost.

\smallskip

Reactions within the Circle are divided:
\begin{itemize}
    \item \textbf{Chris Everett}, a city boy pragmatic and future-oriented, sees opportunity in the investment.
    \item \textbf{Bernadette} and \textbf{Arnie}, uncompromising in their devotion to mathematical purity, reject the idea outright.
    \item \textbf{Emmi} remains torn, fascinated by the interdisciplinary potential but wary of external control.
\end{itemize}

\subsection*{3. Willard’s Hedge Fund Network}

Just as there is an air of inevitability that this is the funding direction they will be taking, an unexpected alternative emerges. \textbf{Willard Jones}, an eccentric but well-connected mathematician, proposes a second path.

\smallskip

His hedge fund contacts manage a \textit{"Green Ethical Tranche"}—a portfolio where high-net-worth individuals invest in fundamental knowledge, not for profit, but as a <<legacy project>>. Unlike Joe’s biomedical investors, they do not seek immediate returns. Instead, they fund intellectual progress, much as others chase glory by building rockets to Mars.

\subsection*{4. The Ethical Investment Pitch}

Willard arranges a meeting with his investors, individuals who have quietly supported \textit{number theory, quantum mechanics, and pure mathematics} in institutions around the world. Their motivations appear ideological rather than financial.

\smallskip

The proposal: The \textbf{Merriweather Circle becomes a self-sustaining institute}, free from commercial entanglements. However, accepting this funding comes with philosophical oversight—these investors have their own grand ideas about the \textit{direction} of human knowledge.

\subsection*{5. The Decision Point}

The Circle now stands at a crossroads:
\begin{enumerate}
    \item Accept the <<big pharma investment>>, securing financial stability but at the cost of potential corporate influence.
    \item Take the <<ethical hedge fund>> backing, retaining independence but facing oversight from a different kind of power.
    \item Reject both offers and risk the <<collapse of Merriweather>>, choosing intellectual purity at the expense of survival.
\end{enumerate}

As the Circle debates, alliances form, betrayals simmer, and a final decision looms while all the while the integrity of the information age shows signs of fraying. In a world where mathematics intersects with financial clout, does it matter of purity survives with compromise?
\newpage
%%%%%%%%%%%%%%%%%%%%%%%%%%%%%
Eleanor Price sat at her favorite desk in the grand reading room of the library, surrounded by towering bookshelves stacked with the accumulated import of generations. The heavy wooden desk was scarred with the marks of long deliberation. Scratches, dents, and worn patches left by scholars before her, each one similarly tortured, wrestling to reveal the most only persistent of truths. The dim golden light from the overhead lamps cast long shadows over her notes, as she scribbled in rumination. Here, in this restless sanctuary, she drafted what would become the foundation of \textbf{Substantive-Mathematics}-a vision of mathematics not as an exploration of abstract islands of truth, but as the boat building and hydrography required to connect those islands.

Mathematics, as is often practiced, is cartography: a formal craft, a mapping of the delineation of abstract landscapes, precise and yet detached. It draws clear borders, marks the terrain, inscribes known structures. But the deeper mathematical endeavor is hydrographic, immersive rather than distant, measuring the shifting depths, tracing currents unseen, mapping not what is fixed but what moves. Substantive mathematics does not seek to merely sketch the outline of knowledge; it sails into the unknown, adapting to the shifting contours of what is \textbf{known}, discovering structure in flux, and refining its instruments with every new crossing. Eleanor outlined her belief systems.

\subsection*{1. Mathematics: A Substrate, Not a Structure}
Mathematics is not Platonic idealism. Not merely a discovery game built on axioms. Not built on unquestionable first principles. It is a \textbf{substrate} enacted by reality. It's axiomatic formalism is mostly illusory. Proof is not as Hilbert might tell you its essence. Rather \textbf{Empiricism}  is.

\subsection*{2. Theorems May Not be so Immutable}
Mathematical truths need not be eternal. History shows this: Euclidean space fell, absolute time collapsed, determinism fractured. Mathematics evolves. Yes reality (should) constrain it. Laws reveal; themselves as a product of their time, shift, adapt \textbf{not fixed.}

\subsection*{3. Mathematics is Empirical}
We reject confinement to rigid formalism rather seeking inspiration from the players who deigned to dabble with life's realities:
\begin{itemize}
\item \textbf{Finance}: Risk, growth, stochastic calculus.\item \textbf{Biology}: Metabolic scaling, self-organization.\item \textbf{Chemical Engineering}: Reaction networks, algebraic mirroring.\item \textbf{Cosmology}: Spacetime structure, quantum symmetries.\item \textbf{Differential Geometry}: Connection, bundle theory and exact sequences.\end{itemize}
Not disparate. Unified expressions of a deeper mathematical fabric.

\subsection*{4. Analogy is Not Just a Tool---It is the Framework}
Mathematics thrives on \textbf{analogy} pointing to a deeper unity.\begin{itemize}
\item Prime numbers : discriminants, determinants, orientation\item Symmetry groups : algebra, conserved quantities\item Financial options : risk, Markov chains, probability\end{itemize}

\subsection*{5. Teleology in Mathematics}
Mathematics possesses tendency and is structure-seeking: Nature follows least action as the universe selects optimal paths. Godel has revealed to mathematician that they need not be so shackled to the axiomatic method: just as valid is to \textbf{discover} through experience, empiricism, and relentless \textbf{pattern-seeking.}

\subsection*{6. The Work Ahead}
We do not carve truths from stone. We unearth them as they are embedded in reality. Connected to science.

\bigskip
\textbf{Eleanor Price}  \textit{Founder, Merriweather Circle}

%%%%%%%%%%%%%%%%%%%#
\newpage
\section*{The cast of characters}
\paragraph{Application to Polynomial Coefficients:} The triangle captures the set of simultaneous linear equations needed to determine the coefficients of a polynomial given its sequence of numbers. For instance, a polynomial \(P(x)\) of degree \(n\) with sequence \((a_1, a_2, \dots, a_n)\) can have its coefficients expressed through the entries of the Worpitsky Triangle.

For Emmi, the rediscovery of the Worpitsky Triangle was a little revelation. Her work on finite differences had already demonstrated how discrete systems could mirror continuous ones. The triangle’s ability to encode difference relationships in its structure aligned perfectly with her research.

\paragraph{Encoding Differences:} Each entry in the Worpitsky Triangle could represent a higher-order difference, capturing the same recursive relationships Emmi used to describe discrete systems.

\paragraph{Polynomial Reconstruction:} The triangle allowed her to reconstruct polynomials from sequences of numbers by solving linear equations directly tied to the triangle’s structure. For example:
A quadratic sequence \(1, 3, 6\) could be reconstructed as:
\[
P(x) = ax^2 + bx + c,
\]
with coefficients derived using the relationships encoded in the triangle.

\subsection*{Bernadette and the Basel Problem}
Bernadette, ever fascinated by the Bernoulli numbers and their ties to infinite series, saw an unexpected connection between the Worpitsky Triangle and her work on the Basel problem.

\paragraph{Bernoulli Numbers and Polynomial Coefficients:} The Worpitsky Triangle, with its ability to generate coefficients for polynomials, offered a new perspective on how Bernoulli numbers arose from recursive relationships. For example:
\begin{itemize}
    \item The coefficients of the Bernoulli polynomials \(B_n(x)\) could be encoded in a Worpitsky-like structure, with the triangle revealing their interplay with finite differences.
    \item Bernadette speculated that the triangle’s structure could help generalize her lighthouse analogy to higher dimensions, where the coefficients of the polynomial would describe the attenuation of light across a lattice.
\end{itemize}

\subsection*{Leary and Sheaf Theory}
Leary John, always looking to tie abstract structures to geometry, saw the Worpitsky Triangle as a tool for understanding the twisting of sheaves and bundles.

\paragraph{Coefficients as Sections:} The entries of the Worpitsky Triangle could represent sections of a line bundle, with each row describing a different degree of twisting over a base space. The triangle’s recursive structure hinted at how sections might vary in response to curvature or topology.

\paragraph{Applications to Twisted Spaces:} Leary began to explore how the triangle’s entries could adapt when the base space was no longer flat:
\begin{itemize}
    \item Over hyperbolic space, the coefficients might involve hyperbolic functions.
    \item Over spherical space, the coefficients could reflect cyclic symmetries, much like in Pascal’s Triangle.
\end{itemize}

\subsection*{Jovannah and Renormalization}
Jovannah Mascerino, with her focus on scaling and regularization, found the Worpitsky Triangle invaluable for understanding renormalization in quantum field theory.

\paragraph{Recursive Regularization:} The triangle’s recursive structure mirrored the step-by-step process of renormalization, where divergences were removed layer by layer. For example:
\begin{itemize}
    \item The entries of the triangle could represent corrections to physical quantities, with each row encoding a new order of approximation.
    \item Jovannah explored how the triangle’s entries scaled with the row number, connecting this behavior to the scaling laws of quantum electrodynamics.
\end{itemize}

As the researchers shared their perspectives, the Worpitsky Triangle became a unifying framework for their work:
\begin{itemize}
    \item \textbf{Emmi’s Finite Differences:} The triangle encoded discrete structures and recursive relationships, bridging the gap between the finite and the infinite.
    \item \textbf{Bernadette’s Bernoulli Numbers:} The triangle offered a new way to understand polynomial coefficients and their ties to infinite series.
    \item \textbf{Leary’s Sheaves:} The triangle hinted at deeper geometric structures, with entries representing sections of twisted bundles.
    \item \textbf{Jovannah’s Scaling:} The triangle’s recursive nature mirrored the processes of renormalization, revealing new connections between geometry and physics.
\end{itemize}

The group gathered in the Studio of Equilibrium, their ideas converging around the Worpitsky Triangle. For each of them, the triangle represented more than just a combinatorial tool—it was a bridge between their diverse fields, a symbol of the power of collaboration.
\bigskip

As Eleanor entered the room, she smiled at the sight of her colleagues deep in thought. The library had become exactly what she had envisioned: a crucible for ideas, where the abstract and the concrete met, twisted, and transformed.





